\section{Aeppli classes and holomorphic Courant extensions}
\label{sec:courant}

A real $(1,1)$-form $\omega$ is \emph{pluriclosed} if
$\partial\bar\partial\omega=0$.
Such a form determines both an Aeppli class and the closed form
$2i\partial\omega$ of type $(2,1)$.
We compute the map between these classes and the holomorphic
Courant extension it determines. Positivity is a separate condition.

Throughout this section, $Y$ is the same selected threefold
as in Section~\ref{sec:bott-chern}, with its fixed nonzero
parameter, integral lifts and adapted gluing.
We use its Hodge table, first Fr\"olicher differentials and
Betti numbers from \eqref{hr:eq:hodge-table},
Proposition~\ref{hr:prop:frolicher} and \eqref{hr:eq:betti}.
The argument uses the additive Aeppli calculation independently
of Theorem~\ref{bc:thm:ring}.

\subsection{The two differential maps from Aeppli cohomology}

Define the real Aeppli group by
\begin{equation}\label{cq:eq:real-Aeppli}
 \mathsf A_{\mathbb R}
 =\frac{\{\omega\in A^{1,1}:\overline\omega=\omega,
                         \ \partial\bar\partial\omega=0\}}
        {\{\partial\lambda+\bar\partial\overline\lambda:
                                             \lambda\in A^{0,1}\}}.
\end{equation}
Its complexification is
$\mathsf A_{\mathbb C}=H_A^{1,1}(Y)$.
Indeed, conjugation preserves the complex numerator and denominator.
Averaging a primitive expression with its conjugate identifies
the real part of the complex denominator with the denominator
in \eqref{cq:eq:real-Aeppli}.

The two differentials induce maps
\begin{equation}\label{cq:eq:two-differentials}
 \partial:\mathsf A_{\mathbb C}\longrightarrow
                         H_{\bar\partial}^{2,1}(Y),\qquad
 \bar\partial:\mathsf A_{\mathbb C}\longrightarrow
                         H_\partial^{1,2}(Y).
\end{equation}
For example, $\partial\bar\partial\omega=0$ makes
$\partial\omega$ $\bar\partial$-closed.
Replacing $\omega$ by $\omega+\partial\lambda+\bar\partial\theta$
changes it by
$\partial\bar\partial\theta=-\bar\partial\partial\theta$,
a Dolbeault boundary. Conjugation gives the other map.

\begin{proposition}\label{cq:prop:two-differentials}
The combined map
\begin{equation}\label{cq:eq:differential-isomorphism}
 (\partial,\bar\partial):\mathsf A_{\mathbb C}
      \xrightarrow{\ \sim\ }
       H_{\bar\partial}^{2,1}(Y)\oplus H_\partial^{1,2}(Y)
\end{equation}
is an isomorphism. Restricting $\partial$ to $\mathsf A_{\mathbb R}$
gives an isomorphism of real vector spaces onto the complex line
$H_{\bar\partial}^{2,1}(Y)$.
\end{proposition}

\begin{proof}
Choose a $\bar\partial$-closed $(1,1)$-form $\zeta$ with
$v=[\partial\zeta]_{\bar\partial}\ne0$.
The rank-one first differential from $(1,1)$ to $(2,1)$
supplies this choice. Put $z=[\zeta]_A$.
Then
\[
 (\partial z,\bar\partial z)=(v,0),\qquad
 (\partial\overline z,\bar\partial\overline z)=(0,\overline v).
\]
Thus $z$ and $\overline z$ are independent.
They form a basis because $h_A^{1,1}=2$.
Both target groups are one-dimensional, proving
\eqref{cq:eq:differential-isomorphism}.
The two source lines are the images of ordinary and conjugate
Dolbeault cohomology in Aeppli cohomology.

A real class has the unique form
$cz+\overline c\,\overline z$.
Its image under $\partial$ is $cv$.
This proves the real isomorphism.
\end{proof}

\subsection{The closed holomorphic recipient}

Let $\Omega_Y^{2,\mathrm{cl}}$ be the sheaf of closed
holomorphic two-forms. The holomorphic Poincar\'e lemma gives
an exact sequence of sheaves of complex vector spaces
\begin{equation}\label{cq:eq:Poincare}
 0\longrightarrow\Omega_Y^{2,\mathrm{cl}}
 \longrightarrow\Omega_Y^2
 \xrightarrow{\ d\ }\Omega_Y^3\longrightarrow0.
\end{equation}
The last map is locally surjective because holomorphic
three-forms are closed in complex dimension three.
The Hodge table gives
$H^0(Y,\Omega_Y^3)=H^1(Y,\Omega_Y^3)=0$.
Consequently the forgetful map is an isomorphism
\begin{equation}\label{cq:eq:forgetful}
 \iota:H^1(Y,\Omega_Y^{2,\mathrm{cl}})
       \xrightarrow{\ \sim\ }H^1(Y,\Omega_Y^2)
       =H_{\bar\partial}^{2,1}(Y).
\end{equation}

The closed-holomorphic cohomology has the smooth representative model
\begin{equation}\label{cq:eq:recipient}
 H^1(Y,\Omega_Y^{2,\mathrm{cl}})
 \simeq
 \frac{\{\tau\in A^{3,0}\oplus A^{2,1}:d\tau=0\}}
      {dA^{2,0}}
\end{equation}
\cite[Theorem~3.3]{GFJSCourant}.
In this model, $\iota$ sends $[\tau]$ to
$[\tau^{2,1}]_{\bar\partial}$.
Changing $\tau$ by $dB$, with $B\in A^{2,0}$, changes that
component by $\bar\partial B$.
This describes the comparison \eqref{cq:eq:forgetful} on representatives.

\begin{theorem}\label{cq:thm:class-map}
The differential defines a complex-linear map
\begin{equation}\label{cq:eq:class-map}
 \Phi:\mathsf A_{\mathbb C}\longrightarrow
                   H^1(Y,\Omega_Y^{2,\mathrm{cl}}),\qquad
 [\omega]_A\longmapsto[2i\partial\omega].
\end{equation}
It has complex rank one and kernel $\mathbb C\overline z$,
with the notation of Proposition~\ref{cq:prop:two-differentials}.
Its restriction
\begin{equation}\label{cq:eq:real-class-map}
 \Phi_{\mathbb R}:\mathsf A_{\mathbb R}
          \xrightarrow{\ \sim\ }H^1(Y,\Omega_Y^{2,\mathrm{cl}})
\end{equation}
is an isomorphism of real vector spaces.
\end{theorem}

\begin{proof}
For an Aeppli representative, $2i\partial\omega$ is $d$-closed.
Under the change
$\omega\mapsto\omega+\partial\lambda+\bar\partial\theta$,
where $\theta$ has type $(1,0)$, it changes by
\begin{equation}\label{cq:eq:Aeppli-gauge}
                   2i\partial\bar\partial\theta
                         =d(-2i\partial\theta).
\end{equation}
This is an allowed boundary in \eqref{cq:eq:recipient}, so
$\Phi$ is well-defined and $\iota\Phi=2i\partial$.
Proposition~\ref{cq:prop:two-differentials} and
\eqref{cq:eq:forgetful} prove its complex rank and kernel.

For a real class,
\[
          \iota\Phi(cz+\overline c\,\overline z)=2icv.
\]
In the real source basis $z+\overline z,i(z-\overline z)$
and target basis $v,iv$, its matrix is
\begin{equation}\label{cq:eq:real-matrix}
                     \begin{pmatrix}0&-2\\2&0\end{pmatrix}.
\end{equation}
Its determinant is four. This proves
\eqref{cq:eq:real-class-map}, including the factor $2i$
and the real structure. The complex kernel meets the real
Aeppli plane only at zero.
\end{proof}

\subsection{The Courant operations and changes of representative}

An exact holomorphic Courant algebroid is an extension
\[
 0\longrightarrow T^{*1,0}Y\longrightarrow Q
        \xrightarrow{\ \pi\ }T^{1,0}Y\longrightarrow0
\]
with its compatible nondegenerate symmetric pairing and Dorfman
bracket on holomorphic sections
\cite[Definitions~3.1--3.2]{GFJSCourant}.
Closed holomorphic two-form transition gauges classify these
objects by $H^1(Y,\Omega_Y^{2,\mathrm{cl}})$
\cite[Theorem~3.3]{GFJSCourant}.

For the representative $\tau$ in \eqref{cq:eq:recipient}, the
global smooth bundle
$Q_\tau=T^{1,0}Y\oplus T^{*1,0}Y$ carries the following structure:
\begin{equation}\label{cq:eq:operations}
 \begin{aligned}
 \bar\partial^\tau(X+\xi)
      &=\bar\partial X+\bar\partial\xi-\iota_X\tau^{2,1},\\
 \pi(X+\xi)&=X,\\
 \langle X+\xi,V+\eta\rangle
      &=\tfrac12\bigl(\xi(V)+\eta(X)\bigr),\\
 [X+\xi,V+\eta]_\tau
      &=[X,V]+\partial(\eta(X))+\iota_X\partial\eta
            -\iota_V\partial\xi+\iota_V\iota_X\tau^{3,0}.
 \end{aligned}
\end{equation}
The bracket is on $\bar\partial^\tau$-holomorphic sections
\cite[Definition~3.4]{GFJSCourant}.

For a smooth two-form $B$, write
$e^B(X+\xi)=X+\xi+\iota_XB$.
If $B$ has type $(2,0)$ and $\tau'=\tau+dB$, then
\begin{equation}\label{cq:eq:B-gauge}
                         e^B:Q_{\tau'}\xrightarrow{\ \sim\ }Q_\tau
\end{equation}
preserves the Dolbeault operator, anchor, pairing and bracket.
Its inverse is $e^{-B}$.
The operator and bracket identities are
\[
 \bar\partial^\tau e^B=e^B\bar\partial^{\tau+dB},\qquad
 [e^Bs,e^Bt]_\tau=e^B[s,t]_{\tau+dB}.
\]
The first uses the $(2,1)$ component $\bar\partial B$,
and the second uses $\partial B$ on holomorphic sections.
In particular, the direction in \eqref{cq:eq:B-gauge} is from
$\tau+dB$ to $\tau$.
Equation~\eqref{cq:eq:Aeppli-gauge} gives the corresponding
explicit gauge for a change of Aeppli representative.

For $\tau=2i\partial\omega$, the $(3,0)$ component vanishes.
The holomorphic sections still satisfy the coupled equation
$\bar\partial^\tau(X+\xi)=0$.
Thus the absence of a $(3,0)$ term in the displayed bracket
does not split the holomorphic Courant algebroid.

\subsection{Reduction of the untwisted smooth algebroid}

The untwisted smooth exact Courant algebroid is $TY\oplus T^*Y$
with anchor the tangent projection, the symmetric pairing of
\eqref{cq:eq:operations}, and bracket
\[
 [X+\xi,V+\eta]_0=[X,V]+\mathcal L_X\eta-\iota_Vd\xi.
\]
Here $TY$ denotes the real tangent bundle.
A \emph{lifting} of $T^{0,1}Y$ is an isotropic involutive
subbundle of its complexification whose anchor maps
isomorphically to $T^{0,1}Y$.

\begin{proposition}\label{cq:prop:smooth-reduction}
Every class in $H^1(Y,\Omega_Y^{2,\mathrm{cl}})$ is realized by
a lifting inside the untwisted smooth exact Courant algebroid.
Its real Aeppli class is unique.
\end{proposition}

\begin{proof}
By Theorem~\ref{cq:thm:class-map}, choose a real
pluriclosed $(1,1)$-form $\omega$ representing its inverse image.
Set $d^c=i(\bar\partial-\partial)$ and $H=-d^c\omega$.
The form $H$ is real and closed.
Since $H_{\mathrm{dR}}^3(Y;\mathbb R)=0$, choose a real
two-form $b$ with $db=H$. Thus
\begin{equation}\label{cq:eq:untwisting}
                         0=-d^c\omega-db.
\end{equation}
The subbundle
\begin{equation}\label{cq:eq:lifting}
                    \ell(\omega,b)=e^{b+i\omega}T^{0,1}Y
\end{equation}
is a lifting \cite[Lemma~3.6]{GFJSCourant}.
Its isotropy follows from skew-symmetry of $b+i\omega$.
For involutivity, the $B$-field bracket identity gives the twist
\begin{equation}\label{cq:eq:reduction-twist}
 db+i\,d\omega
     =-i(\bar\partial-\partial)\omega
                         +i(\partial+\bar\partial)\omega
     =2i\partial\omega.
\end{equation}
Contraction with two $(0,1)$ vectors annihilates this $(2,1)$-form.
Since $T^{0,1}Y$ is involutive, so is \eqref{cq:eq:lifting}.

The quotient $\ell^\perp/\ell$ has the smooth identification
\begin{equation}\label{cq:eq:reduction-map}
 \ell^\perp/\ell\longrightarrow Q_{2i\partial\omega},\qquad
                   [e^{b+i\omega}X+\xi]\longmapsto X+\xi.
\end{equation}
The pairing and anchor are unchanged.
For a $(0,1)$ vector field $W$, the mixed bracket with its lift
defines the Dolbeault operator on the quotient:
\[
        \bar\partial^\ell_W[s]=[e^{b+i\omega}W,s]\bmod\ell.
\]
Substituting \eqref{cq:eq:reduction-twist} gives
$\bar\partial^{2i\partial\omega}$ in \eqref{cq:eq:operations}.
The bracket of holomorphic sections gives its displayed Dorfman
bracket by the same identity
\cite[Lemma~3.8]{GFJSCourant}.
Thus \eqref{cq:eq:reduction-map} preserves all the stated operations.
Uniqueness of the real Aeppli class follows from
\eqref{cq:eq:real-class-map}. The representatives $\omega$ and $b$
need not be unique.
\end{proof}

The real primitive $b$ is smooth gauge data.
It is not a holomorphic $(2,0)$ gauge trivializing a nonzero
class in \eqref{cq:eq:recipient}.
The proposition constructs no positive representative of an
Aeppli class.

\subsection{A necessary condition for a pluriclosed metric}

The cusp supplies a direct test for a positive representative.
Its normalized component is the toric del Pezzo surface of degree
six. Its six boundary curves are copies of $\mathbb P^1$.
The free period quotient identifies them in three pairs, as in
Proposition~\ref{c26:prop:quotient}.
The adapted identification of Section~\ref{hol:sec:adapted-fan}
places this cusp in the selected $Y$.
Restricting its normalization map to any boundary curve gives
a nonconstant holomorphic map
$j:\mathbb P^1\to Y$.
Freeness and the local quotient charts ensure nonconstancy.

\begin{corollary}\label{cq:cor:positive-class}
If $Y$ admits a positive pluriclosed $(1,1)$-form $\omega$, then
\[
             [\omega]_A\ne0,\qquad
             [2i\partial\omega]\ne0
                    \quad\text{in }H^1(Y,\Omega_Y^{2,\mathrm{cl}}).
\]
Its reduced holomorphic Courant algebroid has no splitting
as an exact holomorphic Courant algebroid.
\end{corollary}

\begin{proof}
Strict positivity and nonconstancy of $j$ give
\[
                            \int_{\mathbb P^1}j^*\omega>0.
\]
Indeed, the pullback is positive wherever $dj\ne0$, hence on
a nonempty open subset.
If $[\omega]_A=0$ in \eqref{cq:eq:real-Aeppli}, write
$\omega=\partial\lambda+\bar\partial\overline\lambda$.
On $\mathbb P^1$ the $(2,0)$ and $(0,2)$ components vanish, so
\[
                  j^*\omega=d\,j^*(\lambda+\overline\lambda).
\]
Stokes' theorem contradicts the positive integral.
Injectivity of $\Phi_{\mathbb R}$ gives the second assertion.
The classification by \eqref{cq:eq:recipient} identifies the
split exact holomorphic Courant algebroid with the zero class.
\end{proof}

This corollary gives a necessary condition for a metric.
The next section identifies the reduced embedded toric curves
and evaluates their normalized periods as $(-2,3,0)$.
Their nonnegative zero cone gives positive exact currents and
excludes every pluriclosed Hermitian metric on this $Y$.
The class map and the holomorphic Courant algebroids constructed
above remain defined for all Aeppli classes; positivity fails
for every representative satisfying the pluriclosed equation.
The constructions also retain the selected nonzero parameter
and gluing. Parameter descent, a filling at $u=0$, and comparison
with another marked family require their respective geometric maps.
A chiral or physical realization requires further operations
and comparison maps on its own carrier.
